\documentclass[conference]{IEEEtran}
\usepackage{graphicx}
\usepackage{cite}

\title{A Clean Starting Point for Your IEEE Conference Paper}
\author{\IEEEauthorblockN{Your Name}
\IEEEauthorblockA{Your Department \and your University \and you@university.edu}}

\begin{document}
\maketitle

\begin{abstract}
Replace this paragraph with a short summary of the problem, your method and
the main result.
\end{abstract}

\section{Introduction}
State the motivation and your contribution in one paragraph.

\section{Method}
Describe what you did, precisely enough to be reproduced.

\section{Results}
\begin{figure}[t]
  \centering
  % \includegraphics[width=\linewidth]{result.png}
  \caption{Describe the result in one sentence.}
  \label{fig:result}
\end{figure}

\section{Conclusion}
Summarise the contribution and the next step.

\begin{thebibliography}{9}
\bibitem{ref1}
A. Author, Title of the reference, in emph{Proc. Conf.}, 2026.
\end{thebibliography}
\end{document}
